\documentclass[11pt, letterpaper, includehead]{article}

%%%%%%%%%%%%%%%%%%%%% Pre-document %%%%%%%%%%%%%%%%%%%%%
\usepackage{fancyhdr}  % Allow for headers
\usepackage{graphicx}  % Allow for figures 
\usepackage{float}     % Allow for figure inserted in specified location
\usepackage{amsmath}   % Allow for aligned math
\usepackage{array}     % Allow for cell width manipulatio
\usepackage{nicematrix}
\usepackage{amssymb} % Uhhhh what was this????

\setlength{\parindent}{0pt} % Remove auto paragraph indents

% Get rid of those big ass margins
\usepackage[margin=1in]{geometry}

% Table cell formatting
\setlength{\arrayrulewidth}{0.25mm}
\setlength{\tabcolsep}{11pt}
\renewcommand{\arraystretch}{1.2}

\begin{document}

Finding time to fall from lab table

$$\Delta y = V_{0_y}t + \frac{1}{2}a_yt^2$$
$$\Delta y = \frac{1}{2}a_yt^2$$
$$t = \sqrt{\frac{2\Delta y}{a_y}}$$


Find x velocity for tests

$$\Delta x = V_{0_x}t + \frac{1}{2}a_xt^2$$
$$\Delta x = V_{0_x}t$$
$$V_{0_x} = \frac{\Delta x}{t}$$






Finding Launch Angle based on Measured Velocity, height of target + launcher\\

Set maximum height to be average of top pole and and bottom pole, find $V_{0_y}$ and therefore $\theta$,
as $V_{0_y} = V_0\cdot\sin\theta$

$$V_y^2 = V_{0_y}^2 + 2a_y\Delta y$$
$$0 = V_{0_y}^2 + 2a_y\Delta y$$
$$-V_{0_y}^2 = 2a_y\Delta y$$
$$V_{0_y}^2 = -2a_y\Delta y$$
$$V_{0_y} = \sqrt{-2a_y\Delta y}$$
$$V_0\cdot\sin\theta = \sqrt{-2a_y\Delta y}$$
$$\sin\theta = \frac{\sqrt{-2a_y\Delta y}}{V_0}$$
$$\theta = \arcsin\left(\frac{\sqrt{-2a_y\Delta y}}{V_0}\right)$$

Finding time until ball reaches max height

$$V_y = V_{0_y} + a_yt$$
$$0 = V_{0_y} + a_yt$$
$$-V_{0_y} = a_yt$$
$$t = \frac{-V_{0_y}}{a_y}$$

Finding x displacement

$$\Delta x = V_{0_x}t + \frac{1}{2}a_xt$$
$$\Delta x = V_{0_x}t$$


%%%%%%%%%%%%%%%%%%%%% Title Page %%%%%%%%%%%%%%%%%%%%%
\begin{titlepage}
  \begin{center}
    \Huge{\textbf{Lab 4}}\\
    \Huge{Projectile Motion}
    \vfill
    \begin{figure}[H] % H makes the figure insert at the position in the document
      \centering 
      \includegraphics[width=6cm]{../logo.png}
    \end{figure}
    \large{\textbf{your name here}}\\
    \large{Julian Barossi, Liam Gilligan, Stephanie L'Heureux}\\
    \vspace{0.5cm}
    \normalsize
    \today
  \end{center}
\end{titlepage}

%%%%%%%%%%%%%%%%%%%%% TABLE OF CONTENTS %%%%%%%%%%%%%%%%%%%%%
\tableofcontents
\pagebreak % Move to next page

% Add a nice fancy header
\pagestyle{fancy}
\fancyhead{}
\fancyhead[C]{\textbf{Lab 4:} Projectile Motion}

\section{Measuring the launch Speed (Horizontal Projections)} % 1
\subsection{Measurements} % 1.1
\subsubsection{Verticle distance} % 1.1.1
\subsubsection{Horizontal distance} % 1.1.2
\subsection{Calculations} % 1.2
\subsubsection{Time in air} % 1.2.1
\subsubsection{Launch speed} % 1.2.2
\subsection{Average launch speed, standard deviation, and a standard error} % 1.3
\subsection{Best value for initial speed} % 1.4
\section {Measuring the Launch Speed (Vertical Projections)} % 2
\subsection{Max y height reached} % 2.1
\subsection{Launch speed} % 2.2
By plugging in the displacement from the end of the spring
in the launcher to the maximum height into $\Delta y$, we
can calculate the launch speed for that measurement.

$$V_y^2 = V_{0_y}^2 + 2a_y\Delta y$$
$$0 = V_{0_y}^2 + 2a_y\Delta y$$
$$-V_{0_y}^2 = 2a_y\Delta y$$
$$V_{0_y}^2 = -2a_y\Delta y$$
$$V_{0_y} = \sqrt{-2a_y\Delta y}$$
\subsection{Average launch speed, standard deviation, and a standard error} % 2.3
\subsection{Best value for initial speed} % 1.4

\section{Comparing your Results} % 3
\textbf{Which is most precise? Does this mean that it is the most accurate?}\\ 

\textbf{Are the measured vertical and horizontal launch speeds the same or is there a
significant difference between them? Should they be identical in theory? Why or
why not?}\\

\textbf{What value for initial speed should you use for making predictions in part 4, and
why?}

\section{Challenge Shot} % 4 
\subsection{Prediction} % 4.1
\subsubsection{Optimal theta} % 4.1.1
\subsubsection{Time to maximum height} % 4.1.2
\subsubsection{Horizontal distance} % 4.1.3
\subsection{Analysis}  % 4.2
% How accurate was the prediction 

\subsection{Recorded time} % 1.2 
\begin{center}
  \begin{tabular}[H]{| m{2cm} | m{2cm} |}
    \hline
    \textbf{Trial} & \textbf{Time (s)} \\
    \hline
    1              & 0.55              \\
    \hline
    2              & 0.69              \\
    \hline
    3              & 0.69              \\
    \hline
    4              & 0.65              \\
    \hline
    5              & 0.67              \\
    \hline
    6              & 0.58              \\
    \hline
    7              & 0.61              \\
    \hline
    8              & 0.64              \\
    \hline
    9              & 0.63              \\
    \hline
    10             & 0.68              \\
    \hline
    11             & 0.61              \\
    \hline
    12             & 0.63              \\
    \hline
    13             & 0.67              \\
    \hline
    14             & 0.64              \\
    \hline
    15             & 0.65              \\
    \hline
    16             & 0.68              \\
    \hline
    17             & 0.66              \\
    \hline
    18             & 0.66              \\
    \hline
    19             & 0.56              \\
    \hline
    20             & 0.66              \\
    \hline
  \end{tabular}
\end{center}

\end{document}